\documentclass{article}
\usepackage{pathfinder-readable}

\title{Tail Risk at the Last Quote: A Conditional Link Between Two Market Papers}

\begin{document}
\maketitle

\begin{abstract}
Struski and colleagues study language-model traders in a double auction, while the companion local
manuscript studies fixed-batch estimates of conditional value at risk in a different allocation
mechanism. We asked whether the sampling issue in the second paper could change a late trader's
choice between taking a sure profit and waiting for a better one. The peers found a valid reversal
for a constructed auction state: an expected one-sample risk estimate favours waiting, while true
tail risk favours trading now. The continuation behaviour and the fill probability were assumed, and
no such state was shown to occur in the reported runs. The thread therefore paused pending a
matched-state check in the auction simulator.
\end{abstract}

\section{Introduction}

\emph{Competitive Market Behavior of LLMs} puts language-model buyers and sellers in a
double auction \cite{Q}. A buyer may accept the best current selling price, or post a better
buying offer and hope someone accepts it later. Each round has at most 300 chances to act,
so a better offer can expire without a trade. The agents are told to maximise profit. The
paper reports that one model often makes tiny price improvements and leaves profitable
trades unfinished; its markets reach lower efficiency than the best-performing model's markets.

\emph{Fixed-Batch CVaR and Participation in a Quadratic Resource-Allocation Mechanism}
studies a separate allocation game \cite{P}. It measures risk with conditional value at risk,
or CVaR: the average loss in the worst chosen fraction of outcomes. Its main example shows
that optimising a fixed-batch estimate, even in expectation over fresh batches, can lead to
a choice that fares worse under true risk. Its equilibrium and welfare bounds require the
assumptions of its quadratic mechanism. The two papers called ``Q'' and ``P'' inside that
manuscript are other works, not the pair considered here.

The connection pursued was a late auction choice. Accepting a standing offer secures a modest gain;
improving one's own offer may yield a larger gain, but may never be filled. The peers asked whether
fixed-batch CVaR could rank those actions differently from true CVaR. This was a counterfactual
question: the auction paper gives its traders no CVaR objective or sampling oracle.

\section{What was done}

The peers first read the auction rules and the risk paper's distinction between true CVaR and the
\emph{expected} CVaR of an empirical batch. Ada also asked whether the auction paper's means across
ten runs could reveal the downside of market efficiency. They cannot: different sets of run outcomes
can have the same mean and very different worst outcomes. This raised a separate question about
market-wide risk.

Emmy then reduced the trading question to two actions at the end of a round. Crossing the standing
offer gives a sure profit. Posting a noncrossing offer gives a larger profit only if a later trader
accepts it. Emmy derived the risk ranking for this choice and placed it in a state allowed by the
auction rules. Ada checked the seller count, the resulting fill probability, and the surplus
arithmetic against those rules. Their derivation and checks are in \url{../emmy/cross_wait.md} and
\url{../ada/cvar_cda_note.tex}.

The peers next tested the interpretation. They separated an expected batch objective from a decision
based on one realised batch, and checked what information the selected trader can actually see. They
also checked whether a local decision result says anything about total market surplus. These checks
narrowed the result to a conditional example.

\section{Results}

Let crossing yield certain profit \(g>0\). Let a specified continuation after posting a better,
noncrossing offer yield profit \(h>g\) with probability \(p\), and zero otherwise. Write \(f=1-p\)
for the chance of no trade, and let \(\beta\) be the fraction of worst losses used in CVaR. A
trader's risk-adjusted utility is the negative CVaR of its loss, where loss is negative profit.
Directly averaging the worst \(\beta\) fraction gives the following utility for waiting:
\[
U_{\mathrm{wait}}=h\max\left\{0,1-\frac{f}{\beta}\right\}.
\]
Crossing has utility \(g\). With one fresh observation, the \emph{expected empirical} CVaR gives
waiting utility \(ph\), just its expected profit. Thus, if \(f\geq\beta\) and \(ph>g\), the expected
one-sample objective selects waiting while true CVaR selects crossing. This is an action-ranking
result under the stated continuation, not a claim about how the observed language models choose.

The constructed state makes the comparison concrete. At the penultimate iteration, suppose all 22
traders remain untraded. A buyer values a unit at \(3.25\), the best ask is \(3.00\), and the buyer
can instead improve the best bid to \(2.25\). Crossing gives \(g=1/4\); if a seller accepts the new
bid next, the buyer earns \(h=1\). Seven sellers in the auction's value grid have costs at most
\(2.25\). Stipulate that each of them crosses if selected next and that every other selection leaves
this buyer without a trade. Uniform selection then gives \(p=7/22\) and \(f=15/22\). For a
worst-half tail, \(\beta=1/2\), the checked rankings are
\[
U_{\mathrm{wait}}=0<\frac14,
\qquad U_{\mathrm{one\ sample}}=\frac7{22}>\frac14,
\qquad U_{\mathrm{two\ samples}}=\frac{49}{484}<\frac14.
\]
The two-sample value follows because, for this tail size, both observations must be successful for
waiting to have positive empirical risk utility. These are expectations over fresh batches. A single
realised one-sample estimate instead favours waiting after a sampled fill and crossing after a
sampled failure. Emmy derived the batch calculation; Ada independently checked the auction count and
inequalities in the peer files named above.

If the standing ask belongs to a seller with cost \(1\), crossing creates surplus \(9/4\). Under the
same stipulated continuation, waiting creates expected surplus \(49/88\), as Ada checked. This
comparison concerns only that constructed state. Its waiting option has true-risk utility zero, the
same as no trade, so the reversal is not the negative-participation result in the risk paper.

\section{Objections and unresolved points}

The fill probability assumes a policy for every next trader and uses private seller costs. The
auction buyer sees only its own value, the book and histories. The peers therefore used \(7/22\) as
an analyst's counterfactual, not a trader's belief. A practical test needs an oracle or a belief
from public history. No supplied evidence shows that this state occurred, or that the auction agents
optimised CVaR; here the one-sample choice also maximises expected profit.

Two market objections remain. A buyer who crosses can displace a higher-value buyer, so more
crossing need not raise surplus. Separately, the sum of individual CVaRs need not equal CVaR of
aggregate loss. The auction reports ten-run means, which do not reveal downside efficiency: ten
values of \(0.91\) and nine values of \(1\) plus one of \(0.1\) share a mean but have different
worst runs. These invented sets do not reproduce the auction's data.

The quadratic taxes, equilibrium assumptions and value-error bounds of the risk paper do not carry
over to this discrete auction decision. The auction's ten evaluation runs are also different from
the number of observations in a trader's CVaR batch. The ledger's outside sources already cover
risk-based continuous double-auction bidding \cite{Vytelingum2004}, CVaR bidding in a
double-auctioned power market \cite{Panda2019}, and sampling bias in CVaR \cite{Tamar2015}. The
recorded searches do not establish priority for the proposed combination.

\section{Why the thread stopped}

The verifier accepted the ranking reversal but paused the thread because it rests on assumed
continuation behaviour and a fill probability unavailable to the traders. No evidence shows that
relevant states occur or change market outcomes. The smallest restart is to count reachable late
states in the auction simulator and estimate action-specific continuation distributions at matched
states with fixed other-trader policies. Only then could tests of public-history beliefs and market
outcomes support a claim about these agents.

\bibliographystyle{plainurl}
\bibliography{references}

\end{document}
